\chapter{Competências para Trabalhar num Futuro Incerto}
\label{ch:futuro}

Depois de estudar cenários, regressamos a uma decisão mais próxima: o que vale a pena aprender e praticar? Não há uma lista de ferramentas que garanta emprego. Há competências que permitem testar oportunidades, perceber limites e mostrar trabalho verificável.

\section{Começar pela tarefa que conheces}
Escolhe uma atividade cujo resultado consigas avaliar: rever um texto, testar uma função, analisar dados ou explicar um conceito do teu domínio. Constrói uma base sem IA, uma versão assistida e uma forma de comparar. O capítulo~\ref{ch:trabalho} mostrou por que ganhos e perdas dependem da tarefa e da experiência.

Uma melhoria pequena e demonstrada é melhor evidência profissional do que um catálogo de modelos usados. Descreve o problema, a alteração, os critérios, os resultados e o que não foi validado. Inclui falhas e uma alternativa quando a ferramenta está indisponível.

\section{Fundamentos que ajudam a verificar}
Programação, dados, probabilidade, segurança e conhecimento de domínio ajudam a reconhecer erros que uma interface fluente pode esconder. Comunicação ajuda a transformar uma intenção numa tarefa que outras pessoas conseguem rever. Nenhuma destas competências elimina incerteza; torna-a mais observável.

Não precisas de aprofundar tudo ao mesmo tempo. Para quem testa software, por exemplo, uma progressão possível junta requisitos, oráculos, automação e observabilidade. Para quem escreve, junta investigação, estrutura, voz e revisão factual. São percursos de aprendizagem propostos, sem uma promessa de retorno económico.

Conserva alguma prática sem assistência nas competências necessárias para verificar o resultado. Se não consegues explicar a alteração que assinas, o problema não fica resolvido por o código compilar ou a prosa soar bem.

\section{Escolher uma primeira entrega}

Uma primeira entrega pode ser pequena. Escolhe-a pelo trabalho que queres aprender a verificar.

\begin{center}
\small
\begin{tabularx}{\textwidth}{@{}p{.20\textwidth}XX@{}}
\toprule
Percurso & Entrega & O que outra pessoa pode verificar \\
\midrule
Escrita e estudo & Um aviso revisto ou um guião de estudo com questões reservadas. & Os factos do aviso ou a resolução de um exercício novo. \\
Dados e testes & Uma consulta com um erro reproduzível ou um caso de acesso cruzado. & Resultado esperado, dado de teste e execução. \\
Construção & Um assistente pequeno com fontes e um caminho de encaminhamento. & Uma resposta sustentada, uma recusa e o efeito de mudar uma fonte. \\
\bottomrule
\end{tabularx}
\end{center}

A mesma estrutura serve para aprofundar. Em vez de acrescentares logo mais ferramentas, escolhe um caso que a primeira entrega falhe. Faz a alteração e volta a correr os casos antigos e o novo. Um portefólio com esta sequência permite acompanhar a tua decisão.

\section{Preparar um portefólio que outra pessoa consiga repetir}
Um projeto de portefólio deve permitir executar ou inspecionar a parte relevante. Identifica dados sintéticos, dependências, versões e comandos. Distingue demonstração, protótipo e serviço operacional. Se o projeto usa uma chave paga, oferece um percurso offline que explique contratos e limites quando for tecnicamente possível.

Três projetos pequenos podem servir objetivos distintos: uma ferramenta determinística para um cálculo; uma aplicação com documentos e autorização; uma comparação entre arquiteturas. Não precisam de fingir escala de produção. Precisam de mostrar decisões e evidência.

O livro fornece exemplos desses três tipos. As 52 inferências históricas mostram um procedimento com falhas reais. O novo piloto bloqueado mostra a importância da integridade do ambiente. Nenhum dos dois se torna menos útil por os limites estarem escritos.

\section{Atualizar conhecimento sem perseguir todos os anúncios}
Mantém um conjunto reduzido de fontes primárias: documentação das ferramentas usadas, investigação relevante para a tarefa e orientação aplicável ao domínio. Verifica primeiro mudanças que afetem o teu sistema: modelo, API, protocolo, dados, licença e obrigações. Uma nova ferramenta que não altera uma decisão tua pode esperar.

Reserva tempo para executar, não apenas ler. Depois de uma experiência, escreve o que aprendeste e a próxima pergunta concreta. Se uma previsão falhar, conserva o registo e explica porquê. Mudar de opinião com evidência é parte da competência, não uma falha de imagem.

\begin{promptlabbox}{Um plano de três meses com entregas verificáveis}
\textbf{Objetivo e contexto:} preparar um percurso pessoal com tempo e recursos reais. Não fornece garantia de emprego ou rendimento.

\textbf{Procedimento:} escolhe uma competência de domínio, uma de implementação e uma de avaliação. Define uma pequena entrega mensal e o teste que a pode reprovar. Confirma recursos e disponibilidade. Podes pedir ao modelo alternativas, mantendo a decisão e o calendário sob tua revisão.

\textbf{Resultado esperado e validação:} ao fim do primeiro mês existe uma entrega inspecionável, com critérios e limites, e o plano é revisto à luz do tempo efetivamente gasto.

\textbf{Falhas possíveis:} inventar cursos, contar ferramentas como competência ou encher o calendário com projetos que não cabem no tempo disponível.
\end{promptlabbox}

O capítulo~\ref{ch:grande-laboratorio} reúne projetos para começar. Quando tiveres uma primeira entrega, escolhe uma comparação que possa mudar a tua próxima decisão.
